% !TEX root = ../Vorlage_DA.tex

%	########################################################
% 					Diverse Anhänge
%	########################################################


%	--------------------------------------------------------
% 	Überschrift, Inhaltsverzeichnis
%	--------------------------------------------------------
\chapter{\LaTeX}

\LaTeX (gesprochen Latech) ist ein Textsatzsystem das speziell für umfangreiche  und komplexe wissenschaftliche, technische und mathematische Dokumente entwickelt wurde.
Man schreibt Quelltext wie bei einem Programm und übersetzt diesen Quelltext in ein PDF Dokument.
Siehe \cite{bib:latexintro}.


%~~~~~~~~~~~~~~~~~~~~~~~~~~~~~~~~~~~~~~~~~~~~~~~~~~~~~~~~~~~~~~~~~~~~~~~~~~~~~~~~
\section{Die Vorlage}

Die \LaTeX\ Quelltexte dieses Dokuments sind gedacht um als Vorlage verwendet zu werden.
Dazu muss der Inhalt durch die eigene Arbeit ersetzt werden.
\verb+Vorlage_DA.tex+ ist das zentrale Dokument, in dieses werden die einzelnen Kapitel inkludiert (siehe tex Files im Unterordner \verb+chapters+).


%~~~~~~~~~~~~~~~~~~~~~~~~~~~~~~~~~~~~~~~~~~~~~~~~~~~~~~~~~~~~~~~~~~~~~~~~~~~~~~~~
\section{Programme}

Programme (Editor + PDF Compiler) für \LaTeX:

Für Mac: MacTeX
\url{https://tug.org/mactex/}

Für Windows: MiKTeX
\url{http://miktex.org}


%~~~~~~~~~~~~~~~~~~~~~~~~~~~~~~~~~~~~~~~~~~~~~~~~~~~~~~~~~~~~~~~~~~~~~~~~~~~~~~~~
\section{LaTeX Quelltext}

Bei LaTeX wird der Text, dessen Gliederung in die Formatierung in puren Textfiles mit der Endung \verb+.tex+ beschrieben.
Die Zeichenkodierung der Files muss UTF-8 sein.

Der "`LaTeX Compiler"' (Programm mit dem Namen pdflatex) übersetzt diese Files in ein PDF Dokument.

%~~~~~~~~~~~~~~~~~~~~~~~~~~~~~~~~~~~~~~~~~~~~~~~~~~~~~~~~~~~~~~~~~~~~~~~~~~~~~~~~
\section{Gerüst}

Die Grundstruktur eines \LaTeX\ Dokuments:
\begin{Verbatim}[frame=single]
\documentclass[a4paper,10pt,final,oneside]{scrartcl}	

\usepackage{anysize}
\usepackage[utf8]{inputenc}
\usepackage[T1]{fontenc}
\usepackage[ngerman]{babel}
\usepackage[pdftex]{hyperref}
\usepackage{url}
 
\begin{document}

Ab hier kommt der Inhalt

\end{document}
\end{Verbatim}

Hinweis:
Durch die Vorlage ist diese Grund-Struktur bereits vorgegeben.

%~~~~~~~~~~~~~~~~~~~~~~~~~~~~~~~~~~~~~~~~~~~~~~~~~~~~~~~~~~~~~~~~~~~~~~~~~~~~~~~~
\section{Formatierungen}

\begin{minipage}[t]{0.6\linewidth}
\begin{Verbatim}[frame=single]
Mehrere     Leerzeichen werden 
ignoriert
ein einfacher Zeilenumbruch gilt 
als Leerzeichen.

Eine leere Zeile kennzeichnet einen 
neuen Absatz.
\end{Verbatim}
\end{minipage}
\begin{minipage}[t]{0.4\linewidth}
Mehrere     Leerzeichen werden ignoriert
ein einfacher Zeilenumbruch 
gilt als Leerzeichen.

Eine leere Zeile kennzeichnet 
einen neuen Absatz.
\end{minipage}

\vspace{1em}\noindent
Diverse Formatierungen:

\begin{tabular}{l|l}
\texttt{Typewriter} & \lstinline!\texttt{...}! o.\ \lstinline!{\ttfamily ...}!\\ 
\textbf{Fett} & \lstinline!\textbf{...}! o.\ \lstinline!{\bfseries ...}!\\ 
\textit{Italic} & \lstinline!\textit{...}! o.\ \lstinline!{\itshape ...}!\\ 
\textsl{Slanted} & \lstinline!\textsl{...}! o.\ \lstinline!{\slshape ...}!\\ 
\textsc{Kapitälchen} & \lstinline!\textsc{...}! o.\ \lstinline!{\scshape ...}!\\ 
\textmd{Normal} & \lstinline!\textmd{...}! o.\ \lstinline!{\mdseries ...}!\\ 
\emph{Emph} & \lstinline!\emph{...}! o.\ \lstinline!{\em ...}!\\ 
\textsf{Sans Serif} & \lstinline!\textsf{...}! o.\ \lstinline!{\sffamily ...}!\\
\underline{Unterstrichen} & \lstinline!\underline{...}!\\
Größen & \lstinline!\tiny! \lstinline!\scriptsize! \lstinline!\footnotesize! \lstinline!\small! \lstinline!\normalsize! \lstinline!\large! \lstinline!\Large!\\
 & \lstinline!\LARGE! \lstinline!\huge! \lstinline!\Huge!\\ 
Zentriert & \lstinline!\begin{center}...\end{center}!\\
Gesch.\ Leerz. & \lstinline!~!\\
Zeilenumbruch & \lstinline!\\! oder \lstinline!\newline!\\
Absatzumbruch & \lstinline!\par! oder leere Zeile.\\

\end{tabular}


%~~~~~~~~~~~~~~~~~~~~~~~~~~~~~~~~~~~~~~~~~~~~~~~~~~~~~~~~~~~~~~~~~~~~~~~~~~~~~~~~
\section{Überschriften}

\begin{Verbatim}[frame=single]
\chapter{Hauptkapitel}
\section{Kapitel}
\subsection{Unterkapitel}
\subsubsection{Unterunterkapitel}
\end{Verbatim}

Die Kapitelnummerierung und das Inhaltsverzeichnis werden automatisch erstellt.


%~~~~~~~~~~~~~~~~~~~~~~~~~~~~~~~~~~~~~~~~~~~~~~~~~~~~~~~~~~~~~~~~~~~~~~~~~~~~~~~~
\section{Bilder einfügen}

Formate: pdf, jpg und png.
Dateien im Verzeichnis \verb+media/images+ ablegen.

\begin{Verbatim}[frame=single]
In Abbildung \ref{fig:htl01} sieht man das Logo der HTL Braunau.
\begin{figure}[H]
  \centering
  \includegraphics[width=0.3\textwidth]{./media/images/htl_c_cmyk_rein.pdf}
  \caption{Logo der HTL Braunau.}
  \label{fig:htl01}
\end{figure}
\end{Verbatim}

\begin{framed}
In Abbildung \ref{fig:htl01} sieht man das Logo der HTL Braunau.
\begin{figure}[H]
	\centering
	\includegraphics[width=0.3\textwidth]{./media/images/htl_c_cmyk_rein.pdf}
  	\caption{Logo der HTL Braunau.}
  	\label{fig:htl01}
\end{figure}
\end{framed}

Hinweis: Dateipfade mit \verb+"/"+ bilden!

Siehe
\url{http://en.wikibooks.org/wiki/LaTeX/Importing_Graphics}

Der Befehl \lstinline{\label} gibt der Abbildung einen eindeutigen Namen.
Durch \lstinline{\ref} wird dieser Name referenziert, d.h. es wird die automatisch generierte Abbildungsnummer eingefügt (dazu muss das \LaTeX\ Dokument 2-mal erstellt werden!)

Siehe \ref{ref:abbildungen}, Seite \pageref{ref:abbildungen} für Abbildungen die mehrere Bilder enthalten.

%~~~~~~~~~~~~~~~~~~~~~~~~~~~~~~~~~~~~~~~~~~~~~~~~~~~~~~~~~~~~~~~~~~~~~~~~~~~~~~~~
\section{Querverweise}

Mit Hilfe von Querverweisen verweist man auf andere Stellen im Dokument. 
Z.B. in der Form: "`Siehe \ref{ref:abbildungen}"' --- es wird die Kapitelnummer bzw. Abbildungsnummer angegeben.

Es kann auf Abbildungen und auf Überschriften verwiesen werden.
Diese erhalten zuerst mit \lstinline{\label} einen Namen.
\begin{Verbatim}[frame=single]
\section{Kapitelname} \label{ref:meinebezeichnung}
\end{Verbatim}

Möchte man auf diese Elemente verweisen gibt man den Namen im \lstinline{\ref} Befehl an.
\begin{Verbatim}[frame=single]
Eine genaue Beschreibung dieses Themas ist 
in Kapitel \ref{ref:meinebezeichnung} zu finden.
\end{Verbatim}

Ein Doppelpunkt als Teil des Namens ist erlaubt. Auf diese Weise können zum Beispiel Namen von Abbildungen, Überschriften und Listings unterschieden werden (\mbox{\lstinline{fig:}/\lstinline{ref:}/\lstinline{code:}}).

\LaTeX\ macht aus Querverweisen automatisch PDF Links.


%~~~~~~~~~~~~~~~~~~~~~~~~~~~~~~~~~~~~~~~~~~~~~~~~~~~~~~~~~~~~~~~~~~~~~~~~~~~~~~~~
\section{Aufzählungen}

\begin{minipage}[c]{0.3\linewidth}
\begin{Verbatim}[frame=single]
\begin{itemize}
\item Eins
\item Zwei
\item Drei
\end{itemize}
\end{Verbatim}
\end{minipage}
\begin{minipage}[c]{0.5\linewidth}
\begin{itemize}
\item Eins
\item Zwei
\item Drei
\end{itemize}
\end{minipage}

\noindent
\begin{minipage}[c]{0.3\linewidth}
\begin{Verbatim}[frame=single]
\begin{enumerate}
\item Eins
\item Zwei
\item Drei
\end{enumerate}
\end{Verbatim}
\end{minipage}
\begin{minipage}[c]{0.5\linewidth}
\begin{enumerate}
\item Eins
\item Zwei
\item Drei
\end{enumerate}
\end{minipage}

%~~~~~~~~~~~~~~~~~~~~~~~~~~~~~~~~~~~~~~~~~~~~~~~~~~~~~~~~~~~~~~~~~~~~~~~~~~~~~~~~
\section{Mathematische Formeln}

\begin{Verbatim}[frame=single]
Abgesetzte Formel:
\begin{equation*}
\frac{1+x}{1-x} \cdot \sqrt[3]{2} \cdot \binom{n}{k}
\end{equation*}
\end{Verbatim}

Abgesetzte Formel:
\begin{equation*}
\frac{1+x}{1-x} \cdot \sqrt[3]{2} \cdot \binom{n}{k}
\end{equation*}

\begin{Verbatim}[frame=single]
Formel im Textfluss:
$\frac{1+x}{1-x}, \sqrt[3]{2}, \binom{n}{k}$
\end{Verbatim}

Formel im Textfluß: $\frac{1+x}{1-x}, \sqrt[3]{2}, \binom{n}{k}$
\\

%~~~~~~~~~~~~~~~~~~~~~~~~~~~~~~~~~~~~~~~~~~~~~~~~~~~~~~~~~~~~~~~~~~~~~~~~~~~~~~~~
\section{Programm Quelltext}

Die Umgebung lstlisting übernimmt das Formatieren von Programmquelltext.

\begin{minipage}[t]{\linewidth}
\begin{Verbatim}[frame=single]
Listing \ref{code:complex} zeigt die Implementierung eines 
besonders komplexen Algorithmus.
\begin{lstlisting}[
    language=java,
    caption={Komplizierter Quelltext.},
    label=code:complex
]
while(x>0) {
	x--;
	// bla bla
}
\end{lstlisting}
\end{Verbatim}
\end{minipage}

\begin{framed}
Listing \ref{code:complex} zeigt die Implementierung eines 
besonders komplexen Algorithmus.
\begin{lstlisting}[
    language=java,
    caption={Komplizierter Quelltext.},
    label=code:complex
]
while(x>0) {
	x--;
	// bla bla
}
\end{lstlisting}
\end{framed}

Siehe auch \url{https://en.wikibooks.org/wiki/LaTeX/Source_Code_Listings}


%~~~~~~~~~~~~~~~~~~~~~~~~~~~~~~~~~~~~~~~~~~~~~~~~~~~~~~~~~~~~~~~~~~~~~~~~~~~~~~~~
\section{Seitenümbruche}

Ein Seitenumbruch kann mit \lstinline{\pagebreak} erzwungen werden.

Soll etwas als ganzes auf der Seite stehen und nicht umgebrochen werden, so kann dieser Teil in eine \lstinline{minipage} Umgebung eingeschlossen werden.

\begin{Verbatim}[frame=single]
\begin{minipage}{\linewidth}
In diesem Teil findet kein Seitenumbruch statt.
\end{minipage}
\end{Verbatim}

Häufig wird dies bei Programmlistings nötig sein:

\begin{Verbatim}[frame=single]
Das Listing \ref{code:codenopagebreak} ist in eine minipage 
eingebunden. Dies wirkt sich so aus, dass dieses Listing nur 
als ganzes auf der Seite steht, sollte es nicht mehr Platz 
haben wird es komplett auf die folgende Seite gesetzt und 
es entsteht ein Leerraum auf der vorhergehenden Seite.

\begin{minipage}{\linewidth}
\begin{lstlisting}[
	language=java,
	caption={Java Quelltext.},
	label=code:codeexample1
]
public void prepend(Node n) {
	n.next=start;
	start=n;
}
\end{lstlisting}
\end{minipage}
\end{Verbatim}

Das Listing \ref{code:codenopagebreak} ist in eine minipage eingebunden.
Dies wirkt sich so aus, dass dieses Listing nur als ganzes auf der Seite steht, sollte es nicht mehr Platz haben wird es komplett auf die folgende Seite gesetzt und es entsteht ein Leerraum auf der vorhergehenden Seite.

\begin{minipage}{\linewidth}
\begin{lstlisting}[
	language=java,
	caption={Java Quelltext.},
	label=code:codenopagebreak
]
public void prepend(Node n) {
	n.next=start;
	start=n;
}
\end{lstlisting}
\end{minipage}


%~~~~~~~~~~~~~~~~~~~~~~~~~~~~~~~~~~~~~~~~~~~~~~~~~~~~~~~~~~~~~~~~~~~~~~~~~~~~~~~~
\section{Quellen und Literatur}

Alle Quellen werden im Dokument \verb+includes/literatur.bib+ im BibTeX Format definiert.

\url{http://www.bibtex.org/Format/}

Im Text wird durch den Befehl \lstinline{\cite} zitiert.
Bsp.:

\begin{Verbatim}[frame=single]
Siehe \cite{bib:latexintro}.
\end{Verbatim}

\begin{framed}
Siehe \cite{bib:latexintro}.
\end{framed}


BibTeX braucht einen eigenen Übersetzungsvorgang der je nach Editor anders gestartet werden muss.


%~~~~~~~~~~~~~~~~~~~~~~~~~~~~~~~~~~~~~~~~~~~~~~~~~~~~~~~~~~~~~~~~~~~~~~~~~~~~~~~~
\section{Links}

Links können in der Form einer URL eingefügt werden.
Dies kann (in kleinerem Umfang) statt Einträgen im Literaturverzeichnis verwendet werden.

\begin{Verbatim}[frame=single]
\url{http://www.orf.at}
\end{Verbatim}
\begin{framed}
\url{http://www.orf.at}
\end{framed}



%~~~~~~~~~~~~~~~~~~~~~~~~~~~~~~~~~~~~~~~~~~~~~~~~~~~~~~~~~~~~~~~~~~~~~~~~~~~~~~~~
\section{Fußnoten}

\begin{Verbatim}[frame=single]
Eine Fußnote kann man einfach%
\footnote{Hier zum Beispiel}
irgendwo in den Text einfuegen.
\end{Verbatim}

Eine Fußnote kann man einfach%
\footnote{Hier zum Beispiel}
irgendwo in den Text einfuegen.


%~~~~~~~~~~~~~~~~~~~~~~~~~~~~~~~~~~~~~~~~~~~~~~~~~~~~~~~~~~~~~~~~~~~~~~~~~~~~~~~~
\section{Tabellen}

Mit der Umgebung \lstinline{tabular} bzw. erweitert: \lstinline{tabularx}.

Siehe \url{https://en.wikibooks.org/wiki/LaTeX/Tables}

%: TODO Tabellen


